\documentclass[10pt]{article}
\usepackage[T1]{fontenc}
\usepackage[utf8]{inputenc}
\usepackage[french]{babel}
\usepackage{amsmath,amssymb,mathrsfs,mathpazo}
\usepackage{geometry}
\geometry{a4paper,margin=2cm}
\usepackage{graphicx,array,tabularx,tabularray,colortbl,multirow,multicol,booktabs,enumitem,pifont,ulem,fancybox,diagbox,cancel}
\usepackage{pstricks,pstricks-add,pst-plot,pst-tree,pst-node,pst-text,pst-func,pst-math,pst-bspline,pst-3dplot,pst-eucl,pst-solides3d}
\let\pscalculate\undefined
\usepackage{pst-bezier}
\usepackage{tikz,pgfplots,tkz-tab}
\usetikzlibrary{arrows,arrows.meta,patterns,calc,positioning,trees,decorations.markings,shapes.geometric,shapes.misc,shapes.arrows,intersections,decorations.pathmorphing,decorations.pathreplacing,backgrounds,shadows,babel}
\pgfplotsset{compat=1.18}
\usepgfplotslibrary{fillbetween}
\usepackage[np]{numprint}
\usepackage{listings,makecell,pas-tableur,twoopt}
\let\barre\undefined
\usepackage[pstricks]{tablvar}
\definecolor{darkgreen}{RGB}{0,100,0}
\providecommand{\red}{\color{red}}
\providecommand{\blue}{\color{blue}}
\providecommand{\green}{\color{green}}
\providecommand{\black}{\color{black}}
\providecommand{\white}{\color{white}}
\providecommand{\gray}{\color{gray}}
\providecommand{\cyan}{\color{cyan}}
\providecommand{\magenta}{\color{magenta}}
\providecommand{\yellow}{\color{yellow}}
\providecommand{\violet}{\color{violet}}
\providecommand{\orange}{\color{orange}}
\providecommand{\euro}{\texteuro}
\providecommand{\EUR}{\texteuro}
\providecommand{\ieme}{\textsuperscript{e}}
\providecommand{\degres}{\ensuremath{^\circ}}
\providecommand{\degree}{\ensuremath{^\circ}}
\providecommand{\vv}[1]{\overrightarrow{#1}}
\providecommand{\Coord}[1]{#1}
\providecommand{\decosix}{}
\providecommand{\starredbullet}{\ensuremath{\bullet}}
\pagestyle{empty}
\setlength{\parindent}{0pt}
\usepackage{adjustbox}\DeclareUnicodeCharacter{2236}{:}\DeclareUnicodeCharacter{2190}{\ensuremath{\leftarrow}}
\providecommand{\R}{\mathbb{R}}
\providecommand{\N}{\mathbb{N}}
\providecommand{\Z}{\mathbb{Z}}
\providecommand{\Q}{\mathbb{Q}}
\providecommand{\D}{\mathbb{D}}
\providecommand{\C}{\mathbb{C}}
\providecommand{\e}{{\mathrm{e}}}
\providecommand{\vect}[1]{\overrightarrow{#1}}
\providecommand{\vectt}[1]{\overrightarrow{\mathrm{#1}}}
\providecommand{\vvt}[1]{\overrightarrow{\mathrm{#1}}}
\providecommand{\np}[2][]{{#2~#1}}
\providecommand{\vv}[1]{{\overrightarrow{#1}}}
\renewcommand{\d}{{\mathrm{d}}}
\providecommand{\degres}{{^\circ}}
\providecommand{\degre}{{^\circ}}
\providecommand{\blue}{\color{blue}}
\providecommand{\red}{\color{red}}
\providecommand{\enskip}{\quad}
\providecommand{\llbracket}{\lbrack\!\lbrack}
\providecommand{\rrbracket}{\rbrack\!\rbrack}
\providecommand{\euro}{{\text{€}}}
\providecommand{\no}{{\text{n°}}}
\providecommand{\interval}[2]{{[#1;#2]}}
\providecommand{\intervFO}[2]{{[#1;#2[}}
\providecommand{\dotfill}{\ldots}

\providecommand{\ds}{\displaystyle}
\providecommand{\textding}[1]{#1}
\providecommand{\Oij}{\(\mathrm{O};\vec{\imath},\vec{\jmath}\)}
\providecommand{\Oijk}{\(\mathrm{O};\vec{\imath},\vec{\jmath},\vec{k}\)}
\providecommand{\Ouv}{\(\mathrm{O};\vec{u},\vec{v}\)}
\providecommand{\intervFF}[2]{\ensuremath{[#1;#2]}}
\renewcommand{\euro}{\ifmmode\text{\texteuro}\else\texteuro\fi}
\DeclareUnicodeCharacter{00B0}{\ifmmode^{\circ}\else\textdegree\fi}
\DeclareUnicodeCharacter{20AC}{\ifmmode\text{\texteuro}\else\texteuro\fi}
% Énoncés : Éducation nationale. Transcription et corrigés : APMEP, Denis Vergès et contributeurs. Découpage et classement : CoopMaths, Sébastien Lozano, Éric Elter et contributeurs. Adaptation de la présentation : Maths Explorer.
% CC BY-SA : https://coopmaths.fr/www/about/
% https://coopmaths.fr/alea/static/eam/2026/tex/eam_specialite_2026_06_ameriquenord_automatismes_cor.tex
\begin{document}
\noindent\textit{APMEP / CoopMaths. CC BY-SA. Présentation adaptée par Maths Explorer.}\par\medskip
\textbf{Tableau récapitulatif des réponses}

\begin{center}

		\begin{tblr}{hlines,vlines,colspec={c*{9}{c}},cells={c,m},row{1}={font=\bfseries}}
			Question & 1 & 2 & 3 & 4 & 5 & 6 & 7 & 8 & 9 \\
			Réponse  & B & B & D & C & B & C & B & C & B \\
		\end{tblr}
	\end{center}

	\medskip

	\textbf{Justifications}

	\begin{enumerate}
		\item \textbf{Bonne réponse : B.}

		$\dfrac{1}{2}+\dfrac{3}{2}\times 4= \dfrac{1}{2}+\dfrac{12}{2}=\dfrac{13}{2}$.

		\item \textbf{Bonne réponse : B.}

		La partie visible ($\np[km^3]{150}$) représente $10~\%$ du volume total $V$.

		Donc \quad $0,10 \times V = 150 \iff V=\dfrac{150}{0,10}=\np[km^3]{1500}$.

		\item \textbf{Bonne réponse : D.}

		$0,845=1-0,155 = 1 + \dfrac{-15,5}{100}$ : le coefficient multiplicateur $0,845$ est donc associé à une évolution dont le taux est de $-15,5~\%$, c'est donc une baisse de $15,5~\%$.

		\item \textbf{Bonne réponse : C.}

		$A(x)=(x+5)(x+8)$ est un polynome du second degré, qui s'annule en $-5$ et $-8$.

		Il est présenté sous sa forme factorisée, avec un coefficient dominant positif (égal à 1), donc les valeurs prises par ce polynome sont positives, sauf entre les deux racines.

		On pouvait aussi établir le tableau de signe à partir du signe des facteurs de degré 1 :

		\begin{center}
				\begin{tikzpicture}
					\tkzTabInit[lgt=1.5,espcl=1.4]
					{$x$/0.7,$x+8$/0.7,$x+5$/0.7,$A(x)$/0.7}
					{$-\infty$,$-8$,$-5$,$+\infty$}
				\tkzTabLine{,-,z,+,,+,}
				\tkzTabLine{,-,,-,z,+,}
				\tkzTabLine{,+,z,-,z,+,}
			\end{tikzpicture}
		\end{center}

		\item Les lettres du mot SINGE sont $\{$S, I, N, G, E$\}$ : $5$ lettres, dont $2$ voyelles (I et E).

		Puisque le singe choisit une lettre au hasard, on est en situation d'équiprobabilité, chaque lettre a autant de chance qu'une autre d'être choisie.

		$P_{M}(V)=\dfrac{\text{nombre de voyelles dans SINGE}}{\text{nombre de lettres de SINGE}}=\dfrac{2}{5}$.

		\item \textbf{Bonne réponse : C.}

		La droite passe par les points $(0 ; 30)$ et $(3 ; 0)$.

		L'ordonnée à l'origine est donc : \quad $30$.

		Le coefficient directeur est donné par :\quad $\dfrac{0-30}{3-0}=-10$.

		Donc la fonction représentée a bien pour expression :\quad $f(x)=-10x+30$.

		\item \textbf{Bonne réponse : B.}

		 $\begin{aligned}
		 	(x+2)^{2}-(1-x)^{2}&=\left(x^{2}+4x+4\right)-\left(1-2x+x^{2}\right)\\
		 	&=x^{2}+4x+4-1+2x-x^{2}\\
		 	&=6x+3
		 \end{aligned}$

		\item \textbf{Bonne réponse : C.}

		$\forall x \in \R$,\quad $2(x-4)-(2x+1)=2x-8-2x-1=-9$.

		L'égalité \quad $-9=0$ \quad est fausse, et ce pour tout réel $x$, puisqu'elle ne dépend pas de $x$ : l'équation n'admet aucune solution.

		\item \textbf{Bonne réponse : B.}

		$E=\dfrac{2\times 3^{2}}{27\times 2^{3}}
		=\dfrac{2\times 3^2}{3^3\times 2^3}
		=\dfrac{2^0 \times 3^0}{2^2 \times 3^1}
		=\dfrac{1}{4\times 3}
		=\dfrac{1}{12}$.
	\end{enumerate}

\end{document}
